\documentclass[11pt]{article}
\usepackage[margin=1in]{geometry}
\usepackage{amsmath,amssymb,mathrsfs}
\usepackage[T1]{fontenc}
\usepackage{lmodern}
\usepackage{hyperref}

\newcommand{\Sph}{\mathscr S}
\newcommand{\K}{\mathcal K}

\begin{document}

\begin{center}
{\Large Literature-implied partial answer}\\[3pt]
{\large The operator-norm case of the convexification Mazur map question}
\end{center}

\paragraph{Status.}
\textbf{literature\_implied\_answer (partial subcase).}

\paragraph{Source question.}
In Section 4 of Javier Alejandro Chavez-Dominguez,
\emph{Uniform homeomorphisms between spheres of unitarily invariant ideals},
arXiv:2305.07169, Theorem \texttt{thm-homeomorphism-convexification}
proves that for \(3\le p<\infty\), and every \(1\)-symmetric sequence
space \(E\), the map
\[
G_p(x)=u|x|^p,\qquad x=u|x|,
\]
is a uniform homeomorphism from \(\Sph(S_{E^{(p)}})\) onto
\(\Sph(S_E)\).  The following remark says that it is not known whether
the restriction \(p\ge 3\) is actually needed, and that the restriction is
inherited from Jocić's Theorem 3.1.

\paragraph{Supporting theorem.}
In Section 10 of A. B. Aleksandrov and V. V. Peller,
\emph{Operator Holder-Zygmund functions}, arXiv:0907.3049, Theorem
\texttt{fcc} proves that if \(0<\alpha<1\) and \(f\in\Lambda_\alpha(\mathbb R)\),
then for self-adjoint operators \(A,B\) and a bounded operator \(R\),
\[
\|f(A)R-Rf(B)\|
 \le C_\alpha \|f\|_{\Lambda_\alpha}
       \|AR-RB\|^\alpha \|R\|^{1-\alpha}.
\]
This is an operator-norm theorem for quasicommutators.

\paragraph{Identification.}
Fix \(p>1\) and set \(\alpha=1/p\).  The function
\[
f(t)=\operatorname{sign}(t)|t|^{1/p}
\]
belongs to \(\Lambda_{1/p}(\mathbb R)\).  Applying the Aleksandrov--Peller
quasicommutator estimate gives
\[
\|f(A)R-Rf(B)\|
 \le C_p\|AR-RB\|^{1/p}\|R\|^{1-1/p}.
\]
Replacing \(B\) by \(-B\) turns the denominator into the anticommutator:
for \(A,B\ge0\),
\[
\|A^{1/p}R+RB^{1/p}\|
 \le C_p\|AR+RB\|^{1/p}\|R\|^{1-1/p}.
\]
After substituting \(A=x^p\) and \(B=y^p\), this is precisely the
operator-norm version of the missing reverse estimate in Lemma
\texttt{lemma-sum} of arXiv:2305.07169.

The source paper already proves the forward uniform continuity of \(G_p\)
for all \(p\ge1\).  The only obstruction to extending Theorem
\texttt{thm-homeomorphism-convexification} below \(p=3\) is the inverse
uniform-continuity estimate for the self-adjoint/anticommutator reduction.
The estimate above supplies it for the operator norm.

\paragraph{Conclusion.}
For the \(1\)-symmetric sequence space \(E=c_0\), so that \(S_E=\K(\ell_2)\)
with the operator norm and \(E^{(p)}=c_0\), the map \(G_p\) is a uniform
homeomorphism
\[
G_p:\Sph(S_{c_0})\longrightarrow \Sph(S_{c_0})
\]
for every \(1<p<\infty\).  Thus the restriction \(p\ge3\) is not needed in
this operator-norm subcase.

\paragraph{Limitations.}
This packet does not solve the full question for arbitrary unitarily
invariant ideals and \(1<p<3\).  The supporting theorem used here is in
operator norm; it does not provide the corresponding estimate for an
arbitrary symmetric ideal norm \(S_E\).

\paragraph{References.}
\begin{itemize}
\item J. A. Chavez-Dominguez, \emph{Uniform homeomorphisms between spheres
of unitarily invariant ideals}, arXiv:2305.07169.
\item A. B. Aleksandrov and V. V. Peller, \emph{Operator Holder-Zygmund
functions}, arXiv:0907.3049.
\end{itemize}

\end{document}
